% !TeX root = ../../Book4.tex
% 纲要标题：### 2.3 正向与逆向系统
% 纲要位置：计划纲要.md，第 3143 行。

\section{正向与逆向系统}
\label{b4:sec:0203}

递增子模的并允许先把有限个元素放入共同阶段再计算；逆系统则要求全部坐标同时相容。这一区别决定了两种极限的元素公式和后面的正合性质。

% 纲要内容（仅作写作计划，尚非教材正文）：
%
% * 滤过余极限即通常的正向极限
% * 逆极限及其元素描述；左正合性的相容提升证明
% * 模范畴中的显式公式
%

\subsection{滤过余极限与正向极限}
\label{b4:subsec:0203:01}

考虑递增的子模族。它们的并仍是子模，因为任意两个元素总能放入同一个较大的子模，再在那里完成加法。这个简单事实指出了一种重要条件：有限次计算所涉及的数据，应当能够在同一个对象中比较。

\begin{definition}[滤过范畴与正向系统]\label{b4:def:filtered-category}
设 \(J\) 为非空小范畴。若 \(J\) 满足以下两个条件：
\begin{enumerate}
\item 对任意对象 \(i,j\)，存在对象 \(k\) 及态射 \(i\rightarrow k,j\rightarrow k\)；
\item 对任意平行态射 \(u,v:i\rightarrow j\)，存在 \(w:j\rightarrow k\)，使 \(wu=wv\)。
\end{enumerate}
则称 \(J\) 为\term[lvguofanchou]{滤过范畴}[filtered category]。以滤过范畴为指标的余极限称为\term[lvguoyujixian]{滤过余极限}[filtered colimit]。设 \(\mathcal C\) 为范畴，\(I\) 为非空有向偏序集，即任意两个元素有共同上界。给定 \(\mathcal C\) 中的对象 \(M_i\) 及态射 \(f_{ij}:M_i\rightarrow M_j\)（\(i\le j\)）；若满足
\[
f_{ii}=1,\qquad f_{jk}f_{ij}=f_{ik}\quad(i\le j\le k)
\]
则称这些对象与态射组成一个\term[zhengxiangxitong]{正向系统}[direct system]。该系统的余极限也称为\term[zhengxiangjixian]{正向极限}[direct limit]，记为 \(\varinjlim_i M_i\)。
\end{definition}
\symidx{\(\varinjlim_i M_i\)：正向系统的余极限}

偏序范畴中两个给定对象之间至多有一条态射，所以滤过定义中的第二个条件自动成立。对于一般范畴，这个条件不能省略：它用来协调同一数据被两种不同路径送到同一对象时产生的差别。例如，由两个对象与两条不同的非恒等态射组成的范畴
\[
\begin{tikzcd}[column sep=large]
0\arrow[r,shift left=0.6ex,"u"]\arrow[r,shift right=0.6ex,"v"']&1
\end{tikzcd}
\]
满足共同目标条件，因为两个对象都能映到 \(1\)。但从 \(1\) 出发只有恒等态射，因而无法通过后复合使 \(u,v\) 相等；它不满足第二条条件，并不是滤过范畴。

\begin{lemma}[有限数据的共同比较]\label{b4:lem:filtered-finite-data}
设 \(J\) 为滤过范畴。从 \(J\) 中指定有限个对象以及它们之间有限条态射后，存在一个对象 \(k\) 及从这些对象到 \(k\) 的态射，使每条指定态射对应的三角形交换。
\end{lemma}
\begin{proof}
先反复使用共同目标条件，把所有指定对象映到一个对象 \(k_0\)，记已选的映射为 \(a_i:i\rightarrow k_0\)。对一条指定态射 \(\alpha:i\rightarrow j\)，要协调的是从 \(i\) 到 \(k_0\) 的两条态射 \(a_i\) 与 \(a_j\alpha\)：
\[
\begin{tikzcd}[column sep=large,row sep=large]
i\arrow[r,"\alpha"]\arrow[dr,"a_i"']&j\arrow[d,"a_j"]&\\
&k_0\arrow[r,"w"]&k_1.
\end{tikzcd}
\]
其中三角形起初未必交换。滤过的第二条件给出 \(w:k_0\rightarrow k_1\)，使 \(wa_i=wa_j\alpha\)。把所有已选映射都后复合 \(w\)，这一条件便成立，先前已成立的等式也仍然成立。对剩余的有限条态射重复此步即可。没有指定对象时，使用 \(J\) 非空即可。
\end{proof}

\begin{proposition}[滤过余极限的元素判据]\label{b4:prop:filtered-colimit-elements}
设 \(J\) 为滤过小范畴，\(D:J\rightarrow\mathbf{Set}\) 为图表。余极限的每个元素可由某个 \(x\in D(i)\) 表示。两个代表 \(x\in D(i)\)、\(y\in D(j)\) 给出同一元素，当且仅当存在 \(u:i\rightarrow k\)、\(v:j\rightarrow k\)，使 \(D(u)x=D(v)y\)。特别地，若 \(x,y\in D(i)\)，则它们在余极限中相等，当且仅当存在一个态射 \(w:i\rightarrow k\)，使 \(D(w)x=D(w)y\)。
\end{proposition}
\begin{proof}
在不交并上，把有这样一个共同目标的两个代表记为 \(x\sim y\)。反身性与对称性直接成立。为证传递性，设 \(x\in D(i)\)、\(y\in D(j)\)、\(z\in D(r)\)，并有
\[
\begin{gathered}
u:i\rightarrow k,\quad v:j\rightarrow k,\qquad D(u)x=D(v)y,\\
u':j\rightarrow l,\quad v':r\rightarrow l,\qquad D(u')y=D(v')z.
\end{gathered}
\]
先用第一个滤过条件选取 \(a:k\rightarrow t\)、\(b:l\rightarrow t\)。从 \(j\) 到 \(t\) 的两条路径 \(av\) 与 \(bu'\) 未必相等，故再用第二个条件选取 \(w:t\rightarrow t'\)，使 \(wav=wbu'\)。于是
\[
D(wau)x=D(wav)y=D(wbu')y=D(wbv')z.
\]
态射 \(wau:i\rightarrow t'\) 与 \(wbv':r\rightarrow t'\) 给出 \(x\sim z\) 的见证，因此 \(\sim\) 是等价关系。

这个关系包含每个图表关系 \(x\sim D(\alpha)x\)。反过来，若两个代表在共同目标中相等，则它们都由图表关系与那个共同像等价，所以 \(\sim\) 恰好是一般余极限公式中的生成等价关系。最后，当两个代表同属 \(D(i)\) 时，共同目标判据可能先给出两条不同的态射 \(u,v:i\rightarrow k\)。将它们后复合为同一条态射，即得到最后一个结论。
\end{proof}

这个判据有两个经常使用的结果：有限多个元素可以同时选在某个共同对象中；有限多个已经在余极限中成立的等式，可以在继续向后映射一次后同时成立。第二句话需要先取各等式的见证，再用\cref{b4:lem:filtered-finite-data}协调这些见证，不能把不同等式的目标对象当作自动相同。

\begin{example}[递增的并与分母]\label{b4:exm:direct-limit-localization}
若 \(M_i\) 是同一模 \(M\) 的有向递增子模族，则其正向极限为 \(\bigcup_iM_i\)：相容同态族在这个并上逐元定义，交叠处的一致性由共同上界保证。

\ProseParagraphBreak
对素数 \(p\)，整数模的系统
\[
\mathbb Z\xrightarrow{\times p}\mathbb Z
\xrightarrow{\times p}\mathbb Z\xrightarrow{\times p}\cdots
\]
以 \(n=0,1,\ldots\) 编号，其正向极限为 \(\mathbb Z[1/p]\)。第 \(n\) 项的 \(a\) 对应 \(a/p^n\)，该对应与过渡映射相容；两个分数相等，当且仅当它们在某个更大指标处的整数代表相等。为验证余极限的泛性质，给定交换群 \(M\) 及同态族 \(\phi_n:\mathbb Z\rightarrow M\)，满足 \(\phi_{n+1}(pa)=\phi_n(a)\)，定义
\[
\Phi:\mathbb Z[1/p]\longrightarrow M,\qquad
\Phi(a/p^n)=\phi_n(a).
\]
若 \(a/p^n=b/p^m\)，取 \(\ell\ge n,m\)，则 \(p^{\ell-n}a=p^{\ell-m}b\)。由同态族的相容性，
\[
\phi_n(a)=\phi_\ell(p^{\ell-n}a)
=\phi_\ell(p^{\ell-m}b)=\phi_m(b),
\]
故 \(\Phi\) 良定义。把两个分数写到同一分母后，由 \(\phi_\ell\) 的加法性可知 \(\Phi\) 为同态。每个分数都有上述形式，所以与全部结构映射相容的同态只能是 \(\Phi\)，唯一性也成立。
\end{example}

\subsection{逆极限及其元素描述}
\label{b4:subsec:0203:02}

正向系统把信息逐步送到后面；逆向系统则记录同一个对象在不同精度下的投影。逆极限要求所有精度的选择同时相容，因此其元素通常是一整族数据，而不是某一个阶段的一个元素。

\begin{definition}[逆向系统与逆极限]\label{b4:def:inverse-system}
设 \(I\) 为非空有向偏序集，\(\mathcal C\) 为范畴。对象 \(M_i\) 及态射 \(\rho_{ij}:M_j\rightarrow M_i\)（\(i\le j\)）若满足
\(\rho_{ii}=1\)、\(\rho_{ij}\rho_{jk}=\rho_{ik}\)，则称为一个\term[nixiangxitong]{逆向系统}[inverse system]，即图表 \(I^{\mathrm{op}}\rightarrow\mathcal C\)。该图表的极限称为\term[nijixian]{逆极限}[inverse limit]，记为 \(\varprojlim_i M_i\)。
\end{definition}
\symidx{\(\varprojlim_i M_i\)：逆向系统的极限}

在集合、群、交换环或模中，逆极限均可写为相容元组：
\[
\varprojlim_iM_i=
\{\,(m_i)\in\prod_iM_i\mid
\rho_{ij}(m_j)=m_i\;\text{对所有}\;i\le j\,\}.
\]
代数运算逐分量进行。投影映射给出一个锥，而任何其他锥只能把元素送到其全部分量组成的元组，因而上述公式满足泛性质。

\begin{example}[整数的相容剩余类]\label{b4:exm:adic-inverse-limit}
对素数 \(p\)，约化映射组成逆向系统
\[
\cdots\longrightarrow\mathbb Z/p^3\mathbb Z
\longrightarrow\mathbb Z/p^2\mathbb Z
\longrightarrow\mathbb Z/p\mathbb Z.
\]
其逆极限环记为 \(\mathbb Z_p\)。一个元素是一族相容剩余类 \((a_n)\)，并不要求存在一个整数同时代表全部 \(a_n\)。每个整数给出这样的族；若其全部剩余类均为零，则它被所有 \(p^n\) 整除，故为零。因此 \(\mathbb Z\rightarrow\mathbb Z_p\) 为单射。

\ProseParagraphBreak
逐步选择代表可将每族唯一写成数字 \(c_0,c_1,\ldots\)，其中 \(0\le c_i<p\)，且
\[
a_n\equiv c_0+c_1p+\cdots+c_{n-1}p^{n-1}\pmod{p^n}.
\]
先由 \(a_1\) 决定 \(c_0\)，再由相容性逐项决定下一位即可。这里的无穷数字首先表达相容剩余类，不涉及实数意义下的无穷级数。
\end{example}
\symidx{\(\mathbb Z_p\)：\(p\)-进整数环}

\begin{example}[逆极限不保持满射]\label{b4:exm:inverse-limit-not-exact}
对 \(n\ge1\)，考虑短正合列
\[
0\longrightarrow2^n\mathbb Z\longrightarrow\mathbb Z
\longrightarrow\mathbb Z/2^n\mathbb Z\longrightarrow0.
\]
向较小指标的过渡映射依次为包含映射、恒等映射和约化映射，故各列组成相容的逆向系统。取逆极限后得到
\[
0\longrightarrow0\longrightarrow\mathbb Z
\longrightarrow\mathbb Z_2.
\]
给定 \((b_n)\in\mathbb Z_2\)，每个 \(b_n\) 单独都能选择整数代表 \(c_n\)。但中间系统的过渡映射是恒等映射，相容提升必须满足 \(c_1=c_2=\cdots\)，也就是由同一个整数代表全部 \(b_n\)。这样的整数未必存在，所以最后的映射不是满射。事实上，令 \(b_n\) 为 \(3\) 在
\(\mathbb Z/2^n\mathbb Z\) 中的逆元。逆元的唯一性说明 \((b_n)\) 相容。若它来自整数 \(b\)，则 \(3b-1\) 被所有 \(2^n\) 整除，故 \(3b=1\)，矛盾。
\end{example}

每个阶段的满射都能逐个提升元素，却未必能够把这些提升选成彼此相容的一族。上例显示，讨论逆极限的正合性时，不能把正向极限的结论直接反转。核中的元素则不同：它在对应的子模中有唯一前像，这种唯一性会迫使各阶段的前像相容。

\begin{proposition}[逆极限的左正合性]\label{b4:prop:inverse-limits-left-exact}
设 \(R\) 为含幺环，\(I\) 为非空小有向偏序集，\(A,B,C\) 为以 \(I\) 为指标的幺性左 \(R\) 模逆向系统。记它们的过渡同态分别为 \(\rho^A_{ij},\rho^B_{ij},\rho^C_{ij}\)（\(i\le j\)）。若同态族 \(u_i:A_i\rightarrow B_i\)、\(v_i:B_i\rightarrow C_i\) 与过渡同态相容，且各阶段的序列
\[
0\longrightarrow A_i\xrightarrow{u_i}B_i
\xrightarrow{v_i}C_i
\]
正合，则诱导的序列
\[
0\longrightarrow\varprojlim_iA_i
\xrightarrow{\overline u}\varprojlim_iB_i
\xrightarrow{\overline v}\varprojlim_iC_i
\]
也正合。因此逆极限函子是左正合的。
\end{proposition}
\begin{proof}
诱导同态逐分量施加 \(u_i\) 与 \(v_i\)。若相容族 \((a_i)\) 满足 \(u_i(a_i)=0\) 对所有 \(i\) 成立，则各 \(u_i\) 的单射性给出 \(a_i=0\)，所以 \(\overline u\) 为单射。由 \(v_iu_i=0\)，还有 \(\operatorname{im}\overline u\subseteq\ker\overline v\)。

反过来，若相容族 \((b_i)\in\varprojlim_iB_i\) 属于 \(\ker\overline v\)，则 \(v_i(b_i)=0\)。由 \(\operatorname{im}u_i=\ker v_i\)，存在 \(a_i\in A_i\) 使 \(u_i(a_i)=b_i\)；再由 \(u_i\) 单射，这个原像唯一。对 \(i\le j\)，过渡同态的相容性给出
\[
u_i\bigl(\rho^A_{ij}(a_j)\bigr)
=\rho^B_{ij}\bigl(u_j(a_j)\bigr)
=\rho^B_{ij}(b_j)
=b_i=u_i(a_i).
\]
由 \(u_i\) 单射，\(\rho^A_{ij}(a_j)=a_i\)。所以 \((a_i)\) 是 \(A\) 的相容族，且 \(\overline u((a_i))=(b_i)\)，得到中间位置的正合性。
\end{proof}

\subsection{模范畴中的显式公式}
\label{b4:subsec:0203:03}

模的极限通过方程截取直积，余极限则通过关系商去直和。这正是\cref{b4:thm:limits-products-equalizers}的加性形式：两个同态 \(s,t\) 的等化子是 \(\ker(s-t)\)，余等化子是 \(\operatorname{coker}(s-t)\)。下面的 \(\delta\) 是极限构造中两个直积同态之差，\(d\) 则是余极限构造中两个直和同态之差，把应当相等的两项之差置为零。这两种公式都适用于任意小指标范畴，不要求指标有向。

\begin{proposition}[模的极限与余极限公式]\label{b4:prop:module-limit-formulas}
设 \(R\) 为含幺环，\(J\) 为小范畴，\(M:J\rightarrow R\text{-}\mathbf{Mod}\) 为幺性左 \(R\) 模图表，记 \(M_j=M(j)\)。以下 \(\iota_j\) 与 \(\iota_\alpha\) 分别表示以对象和态射为指标的直和的结构包含。定义
\[
\begin{aligned}
\delta:\prod_{j\in\operatorname{Ob}J}M_j
&\longrightarrow\prod_{\alpha:j\rightarrow k}M_k,
&
\delta((m_j))_\alpha&=M(\alpha)m_j-m_k,\\
d:\bigoplus_{\alpha:j\rightarrow k}M_j
&\longrightarrow\bigoplus_{j\in\operatorname{Ob}J}M_j,
&
d(\iota_\alpha(m))&=\iota_k(M(\alpha)m)-\iota_j(m).
\end{aligned}
\]
则 \(\lim_JM=\ker\delta\)，而
\[
\operatorname{colim}_JM
=\operatorname{coker}d
=\left(\bigoplus_jM_j\right)/S,
\]
其中关系子模为
\[
S=\bigl\langle\,\iota_k(M(\alpha)m)-\iota_j(m)\mid \alpha:j\rightarrow k\;\text{在}\;J\;\text{中},\ m\in M_j\,\bigr\rangle_R.
\]
\end{proposition}
\begin{proof}
第一式的核正好由相容元组组成。第二式中，直和的每个元素只有有限支集，所以 \(d\) 良定义；其像是公式所列关系生成的子模。给定一族同态 \(f_j:M_j\rightarrow N\)，它们唯一给出直和上的同态。该同态在 \(\operatorname{im}d\) 上为零，当且仅当 \(f_kM(\alpha)=f_j\) 对所有 \(\alpha\) 成立。这正是余锥条件，因此它唯一下降到所列余核，得到泛性质。
\end{proof}

\begin{proposition}[滤过模图表的单阶段表示]\label{b4:prop:filtered-module-elements}
设 \(R\) 为含幺环，\(J\) 为滤过小范畴，\(M:J\rightarrow R\text{-}\mathbf{Mod}\) 为幺性左 \(R\) 模图表，记 \(M_i=M(i)\)。其模余极限的每个元素均可表示为某个 \(m\in M_i\) 的像。两个代表 \(m\in M_i\)、\(n\in M_j\) 给出同一元素，当且仅当存在 \(u:i\rightarrow k\)、\(v:j\rightarrow k\)，使 \(M(u)m=M(v)n\)。特别地，\(m\in M_i\) 的像为零，当且仅当某个 \(u:i\rightarrow j\) 满足 \(M(u)m=0\)。
\end{proposition}
\begin{proof}
先在底集合的滤过余极限上定义运算。对 \(m\in M_i,n\in M_j\)，选 \(u:i\rightarrow k\)、\(v:j\rightarrow k\)，定义
\[
[m]+[n]=[M(u)m+M(v)n].
\]
若另选 \(u':i\rightarrow l\)、\(v':j\rightarrow l\)，则\cref{b4:lem:filtered-finite-data}给出 \(a:k\rightarrow t\)、\(b:l\rightarrow t\)，使 \(au=bu'\)、\(av=bv'\)。同态的线性性给出
\[
M(a)\bigl(M(u)m+M(v)n\bigr)
=M(b)\bigl(M(u')m+M(v')n\bigr),
\]
所以两个共同目标算出的和表示同一个类。

若更换代表为 \(m'\in M_{i'}\)、\(n'\in M_{j'}\)，且 \([m]=[m']\)、\([n]=[n']\)，元素判据分别给出
\[
\begin{gathered}
p:i\rightarrow k_0,\quad p':i'\rightarrow k_0,\qquad M(p)m=M(p')m',\\
q:j\rightarrow k_1,\quad q':j'\rightarrow k_1,\qquad M(q)n=M(q')n'.
\end{gathered}
\]
选取 \(a:k_0\rightarrow t\)、\(b:k_1\rightarrow t\)，就有
\[
M(ap)m+M(bq)n=M(ap')m'+M(bq')n'.
\]
由已证明的共同目标无关性，可用 \(t\) 计算两组代表的和，因此更换代表也不改变结果。标量乘法定义为 \(r[m]=[rm]\)；代表在共同目标中的相等，经线性映射仍保持乘以 \(r\) 后的相等，故它也良定义。全部模公理只涉及有限多个元素，可送到一个共同对象后由该模中的公理验证。

所得模接收各 \(M_i\) 的相容线性映射；任意模值余锥在底集合层面诱导的唯一映射保持加法与标量乘法，所以也是模同态。因此这个模就是模范畴的余极限，底集合的元素判据随之成立。零元可在任意对象中取零代表，对同阶段的两个元素使用单个后续态射的判据，即得最后一句。
\end{proof}

\begin{proposition}[滤过余极限的正合性]\label{b4:prop:filtered-colimits-exact}
设 \(R\) 为含幺环，\(J\) 为滤过小范畴，\(A,B,C:J\rightarrow R\text{-}\mathbf{Mod}\) 为幺性左 \(R\) 模图表，记 \(A_j=A(j)\)、\(B_j=B(j)\)、\(C_j=C(j)\)。若自然变换 \(A\xrightarrow{u}B\xrightarrow{v}C\) 在每个对象 \(j\) 上给出短正合列
\[
0\longrightarrow A_j\xrightarrow{u_j}B_j
\xrightarrow{v_j}C_j\longrightarrow0,
\]
则诱导的序列
\[
0\longrightarrow\operatorname{colim}_JA
\xrightarrow{\overline u}\operatorname{colim}_JB
\xrightarrow{\overline v}\operatorname{colim}_JC
\longrightarrow0
\]
也是短正合列。
\end{proposition}
\begin{proof}
若 \(a\in A_i\) 的类被 \(\overline u\) 送到零，则存在 \(\alpha:i\rightarrow j\)，使
\(B(\alpha)u_i(a)=0\)。由自然性，这等于
\(u_jA(\alpha)(a)\)。\(u_j\) 为单射，故 \(A(\alpha)(a)=0\)，于是 \(a\) 在余极限中为零，证明了单射性。

若 \(b\in B_i\) 的类属于 \(\ker\overline v\)，则某个 \(\alpha:i\rightarrow j\) 使
\(v_jB(\alpha)(b)=0\)。阶段 \(j\) 的正合性给出 \(a\in A_j\)，满足 \(u_j(a)=B(\alpha)(b)\)，所以 \([b]=\overline u([a])\)。反向包含由 \(v_ju_j=0\) 得到。最后，任意 \(c\in C_i\) 可在同一阶段提升为某个 \(b\in B_i\)，故 \(\overline v\) 为满射。
\end{proof}

这项论证适用于任意过渡同态。关键在于：一个元素在余极限中为零，就能在某个后续对象中变为零。有限多个这样的等式还可放到同一阶段处理，这也将解释滤过余极限为什么能与有限极限交换。
